% Ultrascannet: MobileFCMViT vs. MambaVision on BUSI
% April 25, 2026

% ============================================================================
% STATUS (remediation in progress — see PLAN in repo history / commit log):
%  - Ablation pipeline had a circularity bug (4/8 variants warm-started from
%    an already-converged checkpoint instead of a common baseline) — FIXED in
%    code (ablation_base.yaml + ablation.py), but the 8 ablation runs below
%    have NOT yet been rerun with the fix, so their numbers are stale.
%  - "FCM-only" / "MobileViTv3-only" standalone architectures were never
%    actually trained; the numbers previously shown for them belonged to
%    other ablation runs (full_fcm_scratch / no_fcm_scratch, mislabeled).
%    Both models are now wired into ablation.py (mobilevitv3_only.py's model
%    class also had a bug — it referenced a timm model name that doesn't
%    exist and would have crashed) but have not yet been run.
%  - No held-out test set existed (train/validation only); a proper deduped
%    train/val/test split now exists (scripts/build_busi_split.py) but the
%    test-set numbers below have not yet been produced.
%  - The training curve below was reconstructed from a run that was silently
%    resubmitted multiple times into the same checkpoint directory; the real
%    per-epoch numbers differ and show a late-training collapse not reflected
%    here. Needs a clean single rerun.
%  - See report_tables_generated.tex (scripts/generate_report_tables.py) for
%    tables sourced directly from result JSON — regenerate after reruns.
% ============================================================================

\section{Introduction}
This report presents a comprehensive evaluation of MobileFCMViT and its ablations on the BUSI breast ultrasound dataset, aiming to surpass the MambaVision baseline (91.02\% top-1 accuracy). We detail the architectures, mathematical foundations, experiment setup, and results, with a focus on Experiment 5 (tuning\_05\_augmix).

\section{Architectures}
\subsection{Fuzzy C-Means (FCM) Only}
The FCM-only model consists of a linear projection of the input image, followed by learnable cluster prototypes. Membership vectors are computed using either softmax or FCM-style assignment, and a classifier predicts the class from the membership. No convolutional backbone is used; the model operates directly on the projected image features.

\subsection{MobileViTv3 Only}
This ablation uses the same custom MobileViTv3-S backbone implementation used by the fusion model (not a timm model — timm has no MobileViTv3 registered; an earlier version of this ablation incorrectly referenced a nonexistent timm model name and would not run) with a classification head. The backbone consists of:
\begin{itemize}
  \item Initial ConvLayer (3x3, stride 2)
  \item Five stages (layer\_1 to layer\_5) of InvertedResidual and MobileViTv3Block modules, with increasing channels and transformer blocks
  \item Final 1x1 ConvLayer expansion
  \item Global average pooling and linear classifier
\end{itemize}
No clustering or FCM modules are present.

\subsection{MobileFCMViT (Fusion)}
MobileFCMViT fuses the MobileViTv3-S backbone with a differentiable feature-space soft clustering layer. The architecture includes:
\begin{itemize}
  \item MobileViTv3-S backbone as above
  \item 1x1 Conv projection to reduce feature dimension
  \item Differentiable soft clustering layer with learnable cluster centers (prototypes)
  \item Membership computation via softmax or FCM formula
  \item Fusion of membership maps with features (attention or concat)
  \item Global average pooling and linear classifier
\end{itemize}

\section{Mathematical Formulation}
\subsection{Fuzzy C-Means Membership}
\begin{align*}
U_{ij} &= \frac{1}{\sum_k (d_{ij}/d_{ik})^{2/(m-1)}} \\
U &= \text{softmax}(-d/\tau)
\end{align*}
where $d_{ij}$ is the distance from sample $i$ to cluster $j$, $m$ is the fuzziness exponent, and $\tau$ is the softmax temperature.

\subsection{Loss Function}
\begin{align*}
L &= L_{\text{task}} + \lambda_1 L_{\text{compactness}} + \lambda_2 L_{\text{entropy}}
\end{align*}
where $L_{\text{compactness}}$ pulls features toward cluster centers and $L_{\text{entropy}}$ encourages confident assignments.

\section{Experiment Setup}
All experiments use the BUSI dataset with three classes (benign, malignant, normal). Training uses AdamW, cosine scheduler, and standard augmentations. Experiment 5 (tuning\_05\_augmix) is prioritized for final results. Validation and test splits are consistent across all runs. Ablation studies systematically disable or modify FCM, backbone, or training components.

\section{Results}
\subsection{Experiment 5 (tuning\_05\_augmix): Training, Evaluation, and Testing}
\textbf{Validation Metrics:}
\begin{itemize}
  \item \textbf{Top-1 Accuracy:} 92.3\%
  \item \textbf{Precision (macro):} 0.932
  \item \textbf{Recall (macro):} 0.904
  \item \textbf{F1 (macro):} 0.917
  \item \textbf{ROC AUC (macro):} 0.973
  \item \textbf{PR AUC (macro):} 0.958
\end{itemize}

\textbf{Per-Class Metrics:}
\begin{tabular}{l|c|c|c|c|c|c}
Class & Acc & Prec & Recall & F1 & ROC AUC & PR AUC \\
\hline
Benign & 0.97 & 0.91 & 0.97 & 0.94 & 0.96 & 0.97 \\
Malignant & 0.86 & 0.92 & 0.86 & 0.89 & 0.96 & 0.93 \\
Normal & 0.89 & 0.96 & 0.89 & 0.92 & 1.00 & 0.98 \\
\end{tabular}

\textbf{Confusion Matrix:}
\begin{tabular}{l|ccc}
 & Pred: Benign & Pred: Malignant & Pred: Normal \\
\hline
True: Benign & 84 & 2 & 1 \\
True: Malignant & 6 & 36 & 0 \\
True: Normal & 2 & 1 & 24 \\
\end{tabular}

\textbf{Training Curve:} pending a clean single rerun of tune\_05\_augmix (see STATUS note above) — the previous curve was reconstructed from a checkpoint directory that had been resubmitted into multiple times, which is not a trustworthy training trajectory. Do not cite the old 0--38 epoch table; the real logged best epoch was 65 (92.95\% val top-1), and later epochs (150+) showed a drop to $\sim$76\% under this run's specific collision history that has not yet been reproduced cleanly.

\subsection{Ablation Results, Tuning Runs, and Held-Out Test Set}
The tables below are generated directly from result JSON files by \texttt{scripts/generate\_report\_tables.py} — never hand-edited — so a row's label always matches the run that produced it.

% ==================================================================
% Auto-generated tables. Regenerate with:
%   python3 scripts/generate_report_tables.py
% Do not hand-edit — edits will be overwritten.
% ==================================================================

\subsection*{Ablation Results (Validation Top-1 Accuracy)}
% Auto-generated by scripts/generate_report_tables.py — do not hand-edit.
\begin{tabular}{l|c}
\textbf{Model} & \textbf{Top-1 Accuracy (\%)} \\
\hline
Full FCM + Pretrained & 92.95 \\
Full FCM, No Pretrained (Scratch) & 85.26 \\
Full FCM, No Mixup & 92.95 \\
Full FCM, No Class Weights & 92.95 \\
Full FCM, No Label Smoothing & 92.95 \\
Full FCM, Dropout 0.3 & 92.95 \\
No FCM (MobileFCMViT arch, FCM losses off), Pretrained & 73.72 \\
No FCM (MobileFCMViT arch, FCM losses off), Scratch & 73.08 \\
FCM-Only (no backbone) & -- (not yet run) \\
MobileViTv3-Only (no FCM module) & -- (not yet run) \\
\end{tabular}

\subsection*{Tuning Runs (Validation Top-1 Accuracy)}
% Auto-generated by scripts/generate_report_tables.py — do not hand-edit.
\begin{tabular}{l|c}
\textbf{Tuning Run} & \textbf{Top-1 Accuracy (\%)} \\
\hline
tune\_01\_lr\_aug & 92.31 \\
tune\_02\_extend & 92.95 \\
tune\_03\_seed2 & 89.74 \\
tune\_04\_focal & 91.03 \\
tune\_05\_augmix & 92.95 \\
tune\_06\_regdrop (datasets2 — not directly comparable) & 88.32 \\
\end{tabular}

\subsection*{Held-Out Test Set Results}
% Auto-generated by scripts/generate_report_tables.py — do not hand-edit.
% Held-out TEST split — never used for checkpoint selection.
\begin{tabular}{l|c|c|c|c|c|c}
Run & Top-1 & Prec (macro) & Recall (macro) & F1 (macro) & ROC-AUC & PR-AUC \\
\hline
\end{tabular}


\section{Discussion}
MobileFCMViT's best single run (tune\_05\_augmix, 92.95\% val top-1) exceeds the MambaVision baseline (91.02\%, computed on the same validation split). However, across the six tuning configs tried, three (tune\_03\_seed2: 89.74\%, tune\_04\_focal: 91.03\%, tune\_06\_regdrop: 88.32\%, the last on a non-comparable data split) sit at or below the baseline — see the auto-generated tuning table above — so this comparison should be read as "the best config we found beats the baseline," not as a robust, seed-independent margin, until multi-seed variance for the winning config is reported (see \texttt{aggregate\_seeds.py}). The ablation claims in the previous version of this report (that ablations "confirm the importance of both the backbone and the FCM module" and that "FCM-only outperforms MobileViTv3-only") were not actually supported by any completed, correctly-labeled run and are withdrawn pending the rerun described in the STATUS note above.

\section{Conclusion}
MobileFCMViT's best individual run outperforms the MambaVision baseline on the BUSI validation split. A defensible claim about the necessity of the FCM module and about generalization (held-out test set, external BUS-UCLM check) requires completing the reruns and multi-seed evaluation described in the STATUS note before this section can be finalized.
